\documentclass[11pt,reqno]{amsart}

\usepackage{lmodern}
\usepackage[T1]{fontenc}
\usepackage[utf8]{inputenc}
\usepackage[english]{babel}
\usepackage{microtype}
\usepackage{amsmath,amssymb,amsthm,mathtools}
\usepackage{booktabs}
\usepackage{xcolor}
\usepackage{hyperref}

% ---------- Theorems ----------
\newtheorem{theorem}{Theorem}[section]
\newtheorem{proposition}[theorem]{Proposition}
\newtheorem{lemma}[theorem]{Lemma}
\newtheorem{corollary}[theorem]{Corollary}
\newtheorem{conjecture}[theorem]{Conjecture}
\newtheorem{observation}[theorem]{Observation}
\theoremstyle{definition}
\newtheorem{definition}[theorem]{Definition}
\newtheorem{example}[theorem]{Example}
\newtheorem{remark}[theorem]{Remark}
\newtheorem{question}[theorem]{Question}

% ---------- Macros ----------
\newcommand{\setZ}{\mathbb{Z}}
\newcommand{\setQ}{\mathbb{Q}}
\newcommand{\setR}{\mathbb{R}}
\newcommand{\Ehr}{L}
\newcommand{\hvec}{\mathbf{h}^{*}}

\newcommand{\defin}[1]{%
\relax\ifmmode%
\textcolor{blue}{#1}%
\else \textcolor{blue}{\emph{#1}}%
\fi%
}

\title{Ehrhart theory of zigzag (fence) poset order polytopes}
\author{Per Alexandersson}
\address{Department of Mathematics, Stockholm University}
\email{per.alexandersson@math.su.se}
\date{\today}

\begin{document}

\begin{abstract}
We study the Ehrhart polynomials and $h^*$-vectors of order polytopes
of fence (zigzag) posets. Using the fast computation methods
from~\cite{AlexanderssonEhrhart}, we tabulate $h^*$-vectors
for $n \leq 12$ and discover that the Ehrhart evaluations at
fixed argument~$s$ satisfy linear recurrences whose characteristic
polynomials have a closed-form generating function of Chebyshev type.
We connect these polytopes to flagged Schur functions via embeddings
into Gelfand--Tsetlin patterns.
\end{abstract}

\maketitle


% =====================================================================
\section{Introduction}
% =====================================================================

Let $F_n$ denote the \defin{fence (zigzag) poset} on $n$ elements,
with cover relations $1 <_{F_n} 2 >_{F_n} 3 <_{F_n} 4 >_{F_n} \dotsb$.
The \defin{order polytope} $\mathcal{O}(F_n)$~\cite{Stanley1986}
consists of all weakly order-preserving maps $F_n \to [0,1]$.

These polytopes are Gorenstein (the fence is self-dual), so their
$h^*$-vectors are palindromic.  Our computations show they are also
ultra-log-concave and real-rooted for all $n \leq 12$.

\begin{remark}
See also the recent work of Br\"and\'en and
Solus~\cite{BrandenSolus2025zigzag} on $h^*$-vectors of
zigzag order polytopes, which establishes real-rootedness and
connections to Chebyshev polynomials.
\end{remark}


% =====================================================================
\section{\texorpdfstring{$h^*$}{h*}-vectors}\label{sec:hstar}
% =====================================================================

Table~\ref{tab:fenceHstar} lists the $h^*$-vectors of
$\mathcal{O}(F_n)$ for $n = 1, \dotsc, 12$, computed via
the frontier DP with Ehrhart--Macdonald reciprocity
from~\cite{AlexanderssonEhrhart}.
All are palindromic (the fence poset is self-dual, so the order polytope
is Gorenstein) and have positive, ultra-log-concave coefficients.
The $h^*$-polynomial has degree $n-2$ for all $n \geq 2$.

\begin{table}[!ht]
\centering
\small
\caption{$h^*$-vectors of order polytopes of fence posets.}
\label{tab:fenceHstar}
\begin{tabular}{@{}rl@{}}
\toprule
$n$ & $h^*(t)$ \\
\midrule
$1$ & $1$ \\
$2$ & $1$ \\
$3$ & $1 + t$ \\
$4$ & $1 + 3t + t^{2}$ \\
$5$ & $1 + 7t + 7t^{2} + t^{3}$ \\
$6$ & $1 + 14t + 31t^{2} + 14t^{3} + t^{4}$ \\
$7$ & $1 + 26t + 109t^{2} + 109t^{3} + 26t^{4} + t^{5}$ \\
$8$ & $1 + 46t + 334t^{2} + 623t^{3} + 334t^{4} + 46t^{5} + t^{6}$ \\
$9$ & $1 + 79t + 937t^{2} + 2951t^{3} + 2951t^{4} + 937t^{5} + 79t^{6} + t^{7}$ \\
$10$ & $1 + 133t + 2475t^{2} + 12331t^{3} + 20641t^{4} + 12331t^{5} + 2475t^{6} + 133t^{7} + t^{8}$ \\
$11$ & $1 + 221t + 6267t^{2} + 47191t^{3} + 123216t^{4} + 123216t^{5}$ \\
     & $\phantom{1} + 47191t^{6} + 6267t^{7} + 221t^{8} + t^{9}$ \\
$12$ & $1 + 364t + 15393t^{2} + 169416t^{3} + 656683t^{4} + 1019051t^{5}$ \\
     & $\phantom{1} + 656683t^{6} + 169416t^{7} + 15393t^{8} + 364t^{9} + t^{10}$ \\
\bottomrule
\end{tabular}
\end{table}


% =====================================================================
\section{A recurrence for Ehrhart evaluations}\label{sec:recurrence}
% =====================================================================

Let $L_n(s) = \Ehr_{\mathcal{O}(F_n)}(s)$ denote the Ehrhart polynomial
of $\mathcal{O}(F_n)$ evaluated at a non-negative integer~$s$.
Computations up to $n = 25$ reveal the following recurrence.

\begin{observation}\label{obs:fenceRecurrence}
For each fixed $s \geq 2$, the subsequences
$(L_{2m+1}(s))_{m \geq 0}$ and $(L_{2m}(s))_{m \geq 1}$
each satisfy a linear recurrence with constant coefficients
given by the nonzero entries of row~$2s$ of the OEIS
triangle~A109466, $T(n,k) = (-1)^{n-k} \binom{k}{n-k}$.

For example:
\begin{itemize}
\item $s=2$: row~4 gives $(1, -3, 1)$, so
  $L_{n+4}(2) = 3 L_{n+2}(2) - L_n(2)$
  for~$n$ of fixed parity.
  The subsequences $(2, 5, 13, 34, 89, \dotsc)$ and
  $(3, 8, 21, 55, 144, \dotsc)$ are bisections of
  the Fibonacci sequence.
\item $s=3$: row~6 gives $(-1, 6, -5, 1)$, so
  $L_{n+6}(3) = 5 L_{n+4}(3) - 6 L_{n+2}(3) + L_n(3)$.
\item $s=4$: row~8 gives $(1, -10, 15, -7, 1)$, yielding an
  order-4 recurrence (in steps of~2).
\end{itemize}
In general, the recurrence for fixed~$s$ has order~$s$
and its characteristic polynomial $p_s(x)$ is read from row~$2s$
of the triangle.
The interleaved sequence $(L_n(s))_{n \geq 1}$
does not satisfy a recurrence from this triangle.

The generating function for the characteristic polynomials is
\begin{equation}\label{eq:fenceCharGF}
  \sum_{s \geq 0} p_s(x)\, y^s
  = \frac{1+y}{1 + (2-x)y + y^2},
\end{equation}
or equivalently, the $p_s$ satisfy the three-term recurrence
\[
  p_s(x) = (x-2)\, p_{s-1}(x) - p_{s-2}(x),
  \qquad p_0 = 1,\; p_1 = x - 1.
\]
\end{observation}


% =====================================================================
\section{Connection to flagged Schur functions}\label{sec:flagged}
% =====================================================================

\begin{remark}\label{rem:fenceFlaggedSchur}
The order polytope of the fence poset on~$n$ elements
can be realized as a face of a Gelfand--Tsetlin polytope:
the zigzag path in the fence embeds as a path through
consecutive rows of a GT pattern of rectangular shape,
with all entries not on the path frozen by the flag
conditions.  This connects the fence Ehrhart polynomial
to flagged Schur functions and may explain the clean
Chebyshev-type structure of~\eqref{eq:fenceCharGF}.
\end{remark}


\begin{thebibliography}{99}

\bibitem{AlexanderssonEhrhart}
P.~Alexandersson,
\emph{Fast computation of Ehrhart polynomials of
Gelfand--Tsetlin polytopes via Macdonald reciprocity},
preprint, 2025.

\bibitem{BrandenSolus2025zigzag}
P.~Br\"and\'en and L.~Solus,
\emph{Symmetric decompositions, triangulations and
real-rootedness},
preprint, 2025.
\url{https://arxiv.org/abs/2503.16403}

\bibitem{Stanley1986}
R.~P.~Stanley,
\emph{Two poset polytopes},
Discrete Comput.\ Geom.\ \textbf{1} (1986), 9--23.

\end{thebibliography}

\end{document}
